\section{Noise- and Fault-Aware Training}\label{sec:training}

Mapping leaves the weights unchanged. Training instead changes them so that the model tolerates the
errors it will meet.

\subsection{Training on Simulated Chips}
Starting from the fine-tuned SVD model of \cref{sec:placement}, every training step draws a new chip.
Each factor of every layer is programmed with the crossbar model of \cref{sec:model}, using the same code
as the evaluation, and the forward pass uses the stored values $\hat W$ in place of the clean weights.
The chip model is not differentiable, so gradients pass through it unchanged, a straight-through
estimator~\cite{bengio2013ste}:
\begin{equation}
W_{\text{fwd}}=W+\operatorname{sg}\big(\hat W-W\big),\qquad
\frac{\partial L}{\partial W}=\frac{\partial L}{\partial W_{\text{fwd}}},
\label{eq:ste}
\end{equation}
where $\operatorname{sg}$ stops the gradient. Because a new chip is drawn at every step, the model cannot
adapt to any single fault map. The training chip has programming noise $\sigma=5\%$ and, where stated,
0.5\% stuck cells. Training runs for 2 epochs over the SST-2 training set with AdamW, learning rate
$10^{-5}$ with linear decay and batch size 32. Dropout stays active as in normal fine-tuning.

\subsection{Variants}
The variants differ only in the chip simulated during training. Every trained model is then tested on
the same hybrid chip, with 20\% of the components in SLC.
\begin{itemize}
\item \textbf{Noise only}: every component on SLC cells, programming noise only. This is plain
noise-aware training.
\item \textbf{Noise + faults}: every component on SLC cells, with noise and stuck cells.
\item \textbf{Tier-aware}: the chip the model will run on. The top 20\% of components of each layer by
$|\sigma_r|$ use SLC cells and the rest MLC cells, with noise and stuck cells. Because singular values
change during training, the SLC set is recomputed from the current values every 200 steps.
\end{itemize}
We call a variant \emph{tier-aware} when it knows the tier, SLC or MLC, in which each component will be
stored, and trains every component under the errors of its own tier. The first two variants are
tier-unaware: every component sees the same cells, as in standard hardware-aware
training~\cite{rasch2023hwa}. The tier-aware variant applies MLC-level errors exactly where MLC will be
used, so the components in the fragile tier learn to tolerate its errors. Comparing the variants in turn
isolates the effect of injecting stuck cells and of knowing the tiers.

The tier-unaware variants use SLC cells for a practical reason. A first attempt placed every component on
MLC cells, the tier-unaware choice that matches the cheaper tier. At 5\% noise and 0.5\% stuck cells such
a chip is at chance accuracy, and the training loss stayed at about 0.70, close to $\ln 2$: the model
learned to hedge between the two classes rather than to tolerate the errors. On SLC cells the loss fell
well below this level from the first steps.
